problem:  
Let \((M^n,g)\) be a compact Riemannian manifold with \(\operatorname{Ric}_g \ge (n-1)g\). For each \(p>1\), denote by \(\lambda_{1,p}(M,g)\) the first nonzero eigenvalue of the \(p\)-Laplacian \(\Delta_p u = \operatorname{div}(|\nabla u|^{p-2}\nabla u)\), and let \(\lambda_{1,p}(S^n)\) be the corresponding eigenvalue on the round unit sphere. It is known that  
\[
\lambda_{1,p}(M,g) \le \lambda_{1,p}(S^n)
\]
by general comparison results (with equality at \(p=2\) characterizing spherical space forms).  

**Open problem:**  
Assume that for some \(p>1\),  
\[
\frac{\lambda_{1,p}(S^n) - \lambda_{1,p}(M,g)}{\lambda_{1,p}(S^n)} \le \varepsilon,
\]
for \(\varepsilon>0\) sufficiently small depending only on \(n\) and \(p\).  

Is it true that there exists a function \(f_{n,p}(\varepsilon)\to 0\) as \(\varepsilon\to 0\) such that \(M\) is \(f_{n,p}(\varepsilon)\)-close (for instance, in the \(C^{1,\alpha}\) or Gromov–Hausdorff sense) to a spherical space form?  

Moreover, does the limiting shape of the extremal \(p\)-eigenfunctions encode the conformal class of the underlying manifold in the small-\(\varepsilon\) regime?  

Why is it a "good" problem:  
This formulation aims at a *nonlinear almost-rigidity phenomenon* for the \(p\)-Laplacian—an analogue of the established Obata rigidity in the linear case \(p=2\), but extended to the nonlinear and potentially degenerate elliptic setting. It naturally bridges nonlinear analysis, geometric measure theory, and global Riemannian geometry by probing the stability of sharp \(p\)-spectral inequalities under curvature constraints.  

If a quantitative stability (or "spectral convergence implies geometric convergence") result could be proved, it would reveal profound structure linking the nonlinear spectral data \(\lambda_{1,p}(M,g)\) to geometric shape, generalizing deep classical insights such as those of Obata, Cheeger–Colding, and others to the nonlinear regime.  

The problem is simple to state, yet touches a frontier area where analytic regularity, metric geometry, and nonlinear PDEs intersect, exemplifying the “narrow entrance, profound depth” nature of a truly good mathematical problem.